\documentclass[11pt]{article}
\usepackage[margin=1in]{geometry}
\usepackage{amsmath,amssymb}
\usepackage{enumitem}
\usepackage{xcolor}
\usepackage{parskip}
\usepackage{booktabs}
\usepackage{array}
\usepackage{hyperref}
\hypersetup{colorlinks=true, urlcolor=blue!45!black, linkcolor=blue!45!black}
\newcommand{\code}[1]{\texttt{\small #1}}
\newcommand{\coordc}{\mathrm{coord}}
\title{\vspace{-2em}\textbf{Implementation Plan --- Spatial Degradation of Hepatocyte Zonation}\\[0.3em]\large Exact inputs, algorithms, formulas, and the code file for every step\\[0.4em]\normalsize\textit{Companion to the stubs in \texttt{src/steps/} and the reference implementation \texttt{src/pipeline.py}}}
\author{}\date{}

\begin{document}
\maketitle
\tableofcontents

\bigskip

This document expands each stub in \code{src/steps/} into exact inputs, algorithms, formulas, data
shapes, outputs, and acceptance checks. It is the ``fill-in-the-blanks'' companion: the docstrings
say \emph{what}; this says \emph{exactly how}. Reference implementations for the statistical steps
already live in \code{src/pipeline.py} (the integrated donor-level script), and each step section
below ends with a \textbf{Code:} line naming the file(s) that implement it.

\section{Conventions used throughout}

\subsection{Paths and the config module}
All paths are imported from \code{src/config.py}; no hard-coded paths appear in the steps.

\begin{itemize}[leftmargin=1.4em]
  \item \code{PAPER1} --- directory of the Paper 1 (disease atlas) inputs:
        \code{counts.mtx}, \code{genes.txt}, \code{barcodes.txt}, \code{cell\_metadata.csv},
        \code{metadata\_all\_cells.csv}.
  \item \code{PAPER2\_TRAIN} --- \code{paper2\_train.npz}, the classifier training bundle.
  \item \code{SIGNATURES} --- the \code{signatures/} directory of zonation gene lists.
  \item \code{FIGURES}, \code{TABLES} --- output directories for artefact figures and CSVs.
  \item \code{config.signature\_files(name)} --- returns the
        \code{(pericentral, periportal)} file pair for a chosen signature set
        (see Section~\ref{sec:sigsets}); \code{config.DEFAULT\_SET} selects the default.
\end{itemize}

\subsection{Matrix orientation, donor, and stage}
\begin{itemize}[leftmargin=1.4em]
  \item \textbf{Matrix.} \code{counts.mtx} is \emph{genes $\times$ cells}
        ($30{,}117 \times 69{,}426$), raw integer UMIs. Row $i \leftrightarrow$
        \code{genes.txt[i]}; column $j \leftrightarrow$ \code{barcodes.txt[j]}.
  \item \textbf{Donor} $=$ \code{metadata\_all\_cells.csv["Patient.ID"]} (about $47$ donors).
  \item \textbf{Stage} $=$ \code{cell\_metadata.csv["stage"]}, an ordered factor mapped to ranks
        $0\ldots4$ by the helper \code{S2R}:
\end{itemize}

\[
\text{Healthy control} < \text{NAFLD} < \text{NASH w/o cirrhosis}
< \text{NASH with cirrhosis} < \text{end stage}
\quad\longmapsto\quad 0,1,2,3,4 .
\]

\subsection{The golden rule --- donor is the unit of inference}
The unit of inference is the \textbf{donor}, not the cell. For every hypothesis we
aggregate cells into \emph{one number per donor}, then run the statistical test on the
$\sim\!47$ donor values. Cell-level p-values are pseudoreplication and are never reported as
the primary result.

\subsection{The A4 table (the shared backbone)}
Step~4 emits a per-cell table consumed by Steps~5--8 with columns

\begin{center}
\code{cell\_id, donor, stage, zonation\_coord, pc, pp, plasticity,}\\
\code{zone\_p0, zone\_p1, zone\_p2, entropy}.
\end{center}

\subsection{Signature sets --- transcriptome-wide by default}
\label{sec:sigsets}
The zonation coordinate is read from a \textbf{gene program}, and the size of that program is a
first-class design choice. The motivation is to read position from the \emph{whole} zonation
program, not from a handful of marker genes.

\paragraph{Primary set: \code{full} (transcriptome-wide).}
The default signature is the full transcriptome-wide zonation program from Paper 2:
\begin{itemize}[leftmargin=1.4em]
  \item \code{signatures/pericentral\_full.txt} --- $1{,}273$ pericentral genes,
  \item \code{signatures/periportal\_full.txt} --- $364$ periportal genes,
\end{itemize}
for a total of $\approx 1{,}637$ \emph{significantly-zonated} hepatocyte genes (Paper 2,
$q < 0.05$, well-expressed, and present in Paper 1). This is \code{config.DEFAULT\_SET = "full"}.
\textbf{Both} the coordinate \emph{and} the classifier use this full set. The classifier's
\code{paper2\_train.npz} cache now stores the \emph{union of all tiers} ($\sim$1{,}700 genes), so any
set can be sliced from it by gene name; re-run \code{src/prep/02\_convert\_paper2\_mat.py} once if the
cache predates this change.

\paragraph{Robustness / baseline alternatives.}
Smaller sets are kept as robustness and baseline checks, each selectable via
\code{config.signature\_files(name)}:
\begin{center}
\begin{tabular}{@{}l l l@{}}
\toprule
\textbf{Set name} & \textbf{Size (PC / PP)} & \textbf{Role} \\
\midrule
\code{full}            & $1273 / 364$ & primary, transcriptome-wide (default) \\
\code{expanded}        & $\sim 102 / 91$ & ranked top genes, $\sim\!100$ each \\
\code{core}            & $13 / 8$ & curated canonical markers \\
\code{paper2\_landmark} & $13 / 8$ exact & Paper 2 landmark baseline \\
\bottomrule
\end{tabular}
\end{center}

\paragraph{The decision: \code{full} is the result, \code{paper2\_landmark} is the audit.}
We report both, but they do \emph{different jobs} --- this is not a 50/50 hedge.
\code{full} is primary because (i)~it \emph{is} the contribution (read position from the entire
program, not markers by eye); (ii)~it is the \emph{more robust} ruler in disease --- the coordinate is
a \emph{mean} of z-scored genes, so disease-corruption of any single marker is diluted (e.g.\
\code{CYP2E1}, canonical pericentral, is induced in NASH for non-spatial reasons and could swing a
13-gene ruler, but is 1-of-$\sim$1{,}600 here); and (iii)~only a large set makes the H2 ``which
programs collapse'' question and its non-circular held-out-gene test (Step~7) meaningful.
\code{paper2\_landmark} is the \emph{credibility control}: a canonical, Paper-2-faithful, fully
interpretable 13$+$8 ruler. If it shows the \emph{same} collapse, the signal is not an artefact of
gene-set size; if the two \emph{disagree}, that is itself a finding (a warning the ruler may be
disease-sensitive). \textbf{Per hypothesis:} H1 (collapse) $+$ ruler/classifier $\to$ \code{full}
primary, \code{paper2\_landmark} control; H2 $\to$ \code{full} with the held-out repeated split;
H3 $\to$ de-zonation score from \code{full}. \code{expanded}/\code{core} are an optional robustness
gradient. \textbf{Why the original plan said ``landmark'':} marker-based reconstruction is the
field-standard, validated method (Paper 2; Halpern et al.\ 2017) --- the conservative starting point;
\code{full} is the deliberate upgrade.

\paragraph{The PC/PP imbalance (1{,}273 vs 364) does not bias the coordinate.}
This is a \emph{feature}-set asymmetry, not ML class imbalance (which concerns label counts, e.g.\ the
classifier's zone terciles). The coordinate is $\operatorname{mean}_z(\text{PC}) -
\operatorname{mean}_z(\text{PP})$; because each arm is a \emph{mean} (not a sum), the 3.5:1 count ratio
carries \emph{no} magnitude bias. It does create a \emph{precision} asymmetry --- the PP arm averages
fewer genes, so it is noisier and the periportal end is estimated less reliably. Mitigations:
(1)~standardise each arm to unit variance across cells \emph{before} subtracting (equalises the arms
regardless of count); (2)~a size-matched run --- top-364 strongest-zonated PC genes vs the 364 PP --- as
a robustness check; (3)~optionally strength-weight genes by Paper~2 zonation effect size, which shrinks
the imbalance by down-weighting the long tail of weak PC genes. The imbalance itself is part biology
(zonation skews pericentral here) and part the selection thresholds ($q<0.05$ $+$ p90 expression $+$
center-of-mass split).

\noindent\textbf{Code:} \code{src/config.py}
(\code{PAPER1}, \code{PAPER2\_TRAIN}, \code{SIGNATURES}, \code{FIGURES}, \code{TABLES},
\code{config.DEFAULT\_SET}, \code{config.signature\_files("full")}).

\section{Step 1 --- Build the ruler (prep)}

\subsection{Inputs}
\begin{itemize}[leftmargin=1.4em]
  \item Paper 2 landmark CSVs $\rightarrow$ \code{signatures/*\_paper2\_landmark.txt}
        (the $20{+}20$ baseline).
  \item \code{supplementary\_table\_8.xlsx}, sheet \code{Hepatocyte} $\rightarrow$
        \code{signatures/*\_expanded.txt} (ranked, $\sim\!100$ each), and the full
        $q<0.05$ program $\rightarrow$ \code{signatures/*\_full.txt}.
  \item \code{combined\_scRNAseq\_atlas\_M5M6M7M8.mat} $\rightarrow$ \code{paper2\_train.npz}.
\end{itemize}

\subsection{Training labels (the classifier's $y$)}
Each Paper 2 nucleus receives a zone label:
\begin{itemize}[leftmargin=1.4em]
  \item \textbf{Ground truth (target).} Port \code{parse\_snRNAseq\_combined\_atlas.m}: Paper 2
        mapped its \emph{spatial} (Visium) zonation onto its snRNA nuclei; use that zone index,
        binned to $3$ classes.
  \item \textbf{Fallback (current \code{paper2\_train.npz}).} Compute on Paper 2's own
        hepatocytes
        \[
          \text{zone\_label} = \operatorname{tercile}\!\big(\text{score(PC)} - \text{score(PP)}\big)
          \in \{0,1,2\}
        \]
        ($0$ portal, $1$ mid, $2$ central).
\end{itemize}

\subsection{Output and acceptance}
\textbf{Artefact A1.} Acceptance: $3$ zone classes present with balanced-ish counts; PC genes
rise and PP genes fall along the zone index (profile plot).

\medskip
\noindent\textbf{Code:} \code{src/prep/02\_convert\_paper2\_mat.py},
\code{src/prep/03\_build\_signatures.py}, the \code{signatures/} directory;
profile plot \code{src/plotting/artefacts.py:plot\_a1\_reference\_profiles}.

\section{Step 2 --- Load \& QC}

\subsection{Inputs}
\code{PAPER1/\{counts.mtx, genes.txt, barcodes.txt, cell\_metadata.csv, metadata\_all\_cells.csv\}}.

\subsection{Algorithm}
\begin{enumerate}[leftmargin=1.6em]
  \item \code{M = scipy.io.mmread(PAPER1/"counts.mtx").tocsc()} $\rightarrow$ genes $\times$ cells
        sparse.
  \item \code{genes = [l.strip() \ldots]}, \code{bars = [l.strip() \ldots]} (order $=$ matrix
        rows / columns).
  \item \code{meta = read\_csv(cell\_metadata).set\_index("cell\_id").reindex(bars)} $\rightarrow$
        \code{stage = meta["stage"]}.
  \item \code{allm = read\_csv(metadata\_all\_cells).set\_index("cell\_id")}; auto-detect the donor
        column from \code{["Patient.ID","patient","donor","orig.ident","sample"]};
        \code{donor = allm[col].reindex(bars)}.
  \item \code{libsize = asarray(M.sum(0)).ravel()} (per-cell total UMIs).
  \item \textbf{QC sanity only (do \emph{not} re-filter).} Confirm the kept cells are
        hepatocytes --- mean \code{ALB} expression is high; print per-stage and per-donor counts.
        Optional sensitivity: flag the lowest-\code{nCount\_RNA} decile for a robustness re-run,
        do not delete.
\end{enumerate}

\subsection{Normalization (used by scoring; the \code{zrows()} helper)}
For a gene row $x$, first CP10k then $\log 1p$:
\begin{align}
  v   &= \log\!\Big(1 + \frac{x}{\text{libsize}} \cdot 10^4\Big), \\
  z   &= \frac{v - \operatorname{mean}(v)}{\operatorname{std}(v)}
         \qquad (\text{if } \operatorname{std}(v)=0,\ z=0).
\end{align}

\subsection{Output and acceptance}
\textbf{Output A2:} \code{M, genes, bars, stage, donor, libsize} plus a per-stage / per-donor
count table.

\textbf{Acceptance:} $69{,}426$ hepatocytes; $5$ stages with counts
$3750 / 7073 / 31903 / 3603 / 23097$; \code{ALB} clearly expressed.

\medskip
\noindent\textbf{Code:} \code{src/steps/step2\_load\_qc.py} (ref \code{src/pipeline.py:load}).

\section{Step 3 --- Harmonize genes}

\subsection{Inputs}
A2 \code{genes} plus the signature gene sets (and \code{paper2\_train.feats} for the classifier).

\subsection{Algorithm}
\begin{enumerate}[leftmargin=1.6em]
  \item Reconcile symbols (and Ensembl IDs where present) across Paper 1, Paper 2, and the
        signatures; report the mapping rate $|\text{shared}| / |\text{signature}|$.
  \item Restrict signatures and classifier features to the shared set. If any signature drops
        below $\sim\!15$ genes, widen it via the \code{expanded} list.
  \item Choose batch-robust features: per-gene z-score (above) and/or a rank transform, so the
        score measures portal-vs-central rather than snRNA-vs-Visium platform.
\end{enumerate}

\subsection{Output and acceptance}
\textbf{Output A3:} a harmonization report (mapping rate, surviving signature sizes) plus the
frozen gene sets. \textbf{Acceptance:} $\ge 80\%$ of signature genes survive the intersection.

\medskip
\noindent\textbf{Code:} \code{src/steps/step3\_harmonize.py}.

\section{Step 4 --- Place every cell}

\subsection{Step 4a --- Signature scoring}
For a gene set $S$, use the \textbf{background-subtracted score} (Scanpy \code{score\_genes}
style). Let $Y$ be the z-scored $\log 1p$-CP10k expression (genes $\times$ cells). For cell $c$,
\[
  \text{score}_S[c] \;=\; \frac{1}{|S|}\sum_{g \in S} Y[g,c]
                       \;-\; \frac{1}{|B|}\sum_{g \in B} Y[g,c],
\]
where $B$ is a control set drawn to match the expression-bin distribution of $S$. Subtracting $B$
removes each cell's overall expression / depth baseline, so the score reflects $S$ specifically.

Then:
\begin{align}
  \text{pc} &= \text{score(PERICENTRAL)}, \qquad \text{pp} = \text{score(PERIPORTAL)}, \\
  \coordc   &= \text{pc} - \text{pp}
               \qquad (\text{signed; } >0 \text{ pericentral, } <0 \text{ periportal}), \\
  \text{plasticity} &= \text{score(PLASTICITY)}
               \qquad (\text{\code{signatures/plasticity.txt}: KRT7, KRT19, SOX9, SOX4, KRT23, NCAM1}).
\end{align}

The default gene sets are the \code{full} program (Section~\ref{sec:sigsets}). The simpler
\code{src/pipeline.py:score} uses the mean of z-scored signature genes \emph{without} an explicit
control set; the \code{score\_genes} background form above is the upgrade.

\medskip
\noindent\textbf{Code:} \code{src/steps/step4\_score.py} (ref \code{src/pipeline.py:score});
plot \code{src/plotting/artefacts.py:plot\_a4\_coordinate}.

\subsection{Step 4b --- Classifier + entropy}
\begin{enumerate}[leftmargin=1.6em]
  \item Train on \code{paper2\_train.npz}: $X$ (nuclei $\times$ feature genes, CP-normalized) and
        $y = \text{zone\_label} \in \{0,1,2\}$. With the \code{full} set the features are the full
        zonated gene set (\emph{not} the $40$ landmarks).
  \item \code{Pipeline(StandardScaler(), LogisticRegression(multi\_class="multinomial", C=\ldots,
        max\_iter=\ldots))}. Prefer a \textbf{calibrated} logistic model (well-behaved
        probabilities) over a random forest.
  \item \textbf{Evaluate on a held-out Paper 2 split FIRST} (\code{train\_test\_split}, stratified
        by donor): accuracy and confusion matrix; must clearly beat chance ($1/3$).
  \item Apply to Paper 1 (same features, same normalization):
        \code{P = clf.predict\_proba(X1)} $\rightarrow$ \code{zone\_p0..2}.
  \item \textbf{Entropy} per cell:
        \[
          H[c] = -\sum_{k} p_k \log p_k
        \]
        ($k$ over the $3$ zones, natural log, range $[0, \log 3]$). Low $H$ $=$ confident zone;
        high $H$ $=$ de-zonated / confused.
  \item \textbf{4c agreement.} \code{spearman(zonation\_coord, expected\_zone)} on Paper 1
        (especially healthy) should exceed $\sim 0.5$.
\end{enumerate}

\textbf{Output A4:} merge \code{zone\_p0..2, entropy} into the per-cell table.

\medskip
\noindent\textbf{Code:} \code{src/steps/step4b\_classifier.py} (ref \code{src/classifier.py}).

\section{Step 5 --- Validate on healthy (GATE)}

\subsection{Input and algorithm}
Input: A4 restricted to Paper 1 healthy donors. For each validation marker,
\[
  \rho = \operatorname{spearman}\big(\coordc[h],\ \text{expr}_{\text{marker}}[h]\big).
\]
\begin{itemize}[leftmargin=1.4em]
  \item Pericentral \code{CYP2E1, CYP1A2, GLUL, PCK2} $\rightarrow$ expect $\rho > 0$.
  \item Periportal \code{ASS1, ALDOB, PCK1, HAL} $\rightarrow$ expect $\rho < 0$.
  \item Classifier entropy in healthy should be low (peaked predictions).
\end{itemize}

\subsection{Output and the gate}
\textbf{Output A5:} a bar of $\rho$ per marker plus the values. \textbf{Acceptance / GATE:}
pericentral $\rho>0$ and periportal $\rho<0$. \emph{Do not start Step 6 until this passes.} If it
fails: widen to \code{expanded} signatures, use rank features, revisit normalization.

\medskip
\noindent\textbf{Code:} \code{src/steps/step5\_validate.py} (ref \code{src/pipeline.py:validate});
plot \code{src/plotting/artefacts.py:plot\_a5\_validation}.

\section{Step 5b --- Is the ruler still a ruler?}

Per stage $s$, on the cells of that stage:
\begin{enumerate}[leftmargin=1.6em]
  \item \textbf{Internal coherence.} Mean pairwise correlation among pericentral genes (and among
        periportal genes). Healthy $=$ high; if it crashes, the module is dissolving (not just
        de-zonating).
  \item \textbf{Cross-program anti-correlation.} $\operatorname{corr}(\text{pc}, \text{pp})$
        across cells. Healthy $\approx$ strongly negative; weakening toward $0$ $=$ true
        de-zonation (robust to global shifts).
  \item \textbf{Split-half reproducibility.} Randomly split each signature in two, build two
        coordinates, and correlate them. A drop means the signature is no longer self-consistent.
  \item \textbf{Program-off vs restriction-lost.} Track each module's mean magnitude (is it
        switched off?) versus the rate of cells co-expressing both ends (is positional restriction
        lost?).
\end{enumerate}

\textbf{Output A5b:} per-stage curves; tells you \emph{which kind} of breakdown each stage is.

\medskip
\noindent\textbf{Code:} \code{src/steps/step5b\_ruler\_diagnostics.py};
plot \code{src/plotting/artefacts.py:plot\_a5b\_ruler}.

\section{Step 6 --- Collapse curve (H1)}

\subsection{Per-donor metrics}
For donors with $\ge 30$ cells, one row per donor:
\begin{align}
  \text{spread}[d]     &= \operatorname{std}\big(\coordc[\text{donor}=d]\big)
                          && \text{(narrows with disease)}, \\
  \text{anticorr}[d]   &= \operatorname{spearman}\big(\text{pc}[\text{donor}=d],\ \text{pp}[\text{donor}=d]\big)
                          && \text{(rises toward } 0), \\
  \text{entropy}[d]    &= \operatorname{mean}\big(H[\text{donor}=d]\big)
                          && \text{(rises, if classifier run)}, \\
  \text{stage\_rank}[d] &= \text{S2R}\big[\operatorname{mode}(\text{stage}[\text{donor}=d])\big].
\end{align}

\subsection{Ordered-trend test on the $\sim\!47$ donor values}
\begin{itemize}[leftmargin=1.4em]
  \item \textbf{Spearman.} \code{rho, p = spearmanr(stage\_rank, metric)}; with no ties
        \[
          \rho = 1 - \frac{6\sum_d d^2}{n(n^2-1)}.
        \]
  \item \textbf{Jonckheere--Terpstra} (ordered alternative across the $5$ stage groups):
        \[
          J = \sum_{a<b} U(\text{group}_a, \text{group}_b),
        \]
        where $U$ counts pairs with $x_b > x_a$. Standardize with the null mean
        \[
          \mu_J = \frac{N^2 - \sum_s n_s^2}{4}
        \]
        and the standard JT variance $\sigma_J^2$ to obtain $z$ and $p$.
\end{itemize}

\subsection{Confidence and negative control}
\begin{itemize}[leftmargin=1.4em]
  \item \textbf{Donor bootstrap CI.} Resample donors with replacement $B = 2000$ times, recompute
        $\rho$, and take the $2.5 / 97.5$ percentiles.
  \item \textbf{Negative control (permutation).} Permute \code{stage\_rank} across donors $B = 2000$
        times, recompute $|\rho|$, and report
        \[
          p_{\text{perm}} = \frac{\#\{\,\text{null} \ge \text{observed}\,\} + 1}{B + 1}.
        \]
\end{itemize}

\textbf{Output A6:} \code{collapse\_per\_donor.csv} plus a plot (donor points $+$ per-stage mean)
and per-metric $\rho$, $95\%$ CI, $p$, $p_{\text{perm}}$.
\textbf{Acceptance:} spread $\downarrow$, anticorr $\to 0$, (entropy $\uparrow$); CI excludes $0$;
$p_{\text{perm}}$ small.

\medskip
\noindent\textbf{Code:} \code{src/steps/step6\_collapse.py} (ref \code{src/pipeline.py:collapse});
plot \code{src/plotting/artefacts.py:plot\_a6\_collapse}.

\section{Step 7 --- Zone-resolved DE (H2)}

\subsection{Bin cells into zones}
By coordinate terciles: \code{terc = quantile(coord, [1/3, 2/3])},
\code{zone = digitize(coord, terc)} $\rightarrow$ $0$ portal, $1$ mid, $2$ central.

\subsection{Pseudobulk (the donor-level move)}
For each zone $z \in \{\text{portal}, \text{central}\}$, for each donor $d$ with $\ge 20$ cells in
that zone, build one profile:
\begin{align}
  \text{cells}  &= (\text{zone}=z)\ \&\ (\text{donor}=d), \\
  \text{cp}     &= \text{per-cell CP-normalized counts} = \text{counts} / \text{colsum}, \\
  \text{PB}[d,:] &= \log\!1p\!\Big(\operatorname{mean}_{\text{cells}}\big(\text{cp}\big) \cdot 10^4\Big)
                    && \text{(donor $\times$ gene)}, \\
  \text{ranks}[d] &= \text{stage\_rank}[d].
\end{align}

\subsection{Test per gene (detected in $\ge 10\%$ of the zone's donors)}
\begin{itemize}[leftmargin=1.4em]
  \item \textbf{MVP.} \code{rho, p = spearman(PB[:,g], ranks)} across donors.
  \item \textbf{Cleaner H2 (the real test).} OLS interaction on the pseudobulk
        \[
          \text{expr}_g \;\sim\; \coordc + \text{stage} + \coordc{:}\text{stage};
        \]
        the $\coordc{:}\text{stage}$ coefficient answers ``does this gene's zonal slope weaken with
        stage?'' --- i.e. lost zonation, not just lost level.
\end{itemize}

\subsection{Multiple testing (BH-FDR)}
Sort the $p$-values ascending; the threshold is $p_{(k)} \le \frac{k}{m}\alpha$; take the largest
such $k$; everything below it is significant; the $q$-value is the monotone-adjusted $p$. Report at
$q < 0.05$.

\subsection{Circularity guards (mandatory)}
The coordinate is built \emph{from} the signature genes, so testing those same genes for
``collapse'' is circular. Two guards:
\begin{itemize}[leftmargin=1.4em]
  \item \textbf{Signature flag.} Flag \code{is\_signature} genes and report significant counts
        BOTH including and EXCLUDING them; ``driver'' claims use non-signature genes.
  \item \textbf{Held-out gene split (the key fix).} Build the coordinate on one random half of the
        gene set and test collapse on the disjoint other half.
\end{itemize}

\paragraph{Repeated random held-out splits.}
Because any single random partition could be lucky, repeat it $K$ times
($K \approx 20$--$50$ random splits) and report the \emph{distribution} of the result across splits:
\begin{itemize}[leftmargin=1.4em]
  \item the \textbf{median effect} across splits, and
  \item the \textbf{fraction of splits} that are significant.
\end{itemize}
A result that holds across most random splits is robust; one that depends on a particular partition
is not. Practical notes:
\begin{itemize}[leftmargin=1.4em]
  \item \textbf{Fix the RNG seed} for reproducibility.
  \item \textbf{Optionally stratify} the split so each half keeps both PC and PP genes.
  \item This is \textbf{meaningful only with the large \code{full} set}, not $20$ landmarks ---
        you cannot split $40$ genes into two informative halves.
\end{itemize}

\textbf{Output A7:} \code{de\_portal.csv}, \code{de\_central.csv}
(\code{gene, effect, p, q, is\_signature}) plus a volcano plot.
\textbf{Acceptance:} the significant set survives excluding signature genes \emph{and} holds across
most random splits.

\medskip
\noindent\textbf{Code:} \code{src/steps/step7\_de.py} (ref \code{src/pipeline.py:de});
plot \code{src/plotting/artefacts.py:plot\_a7\_volcano}.

\section{Step 8 --- De-zonation $\leftrightarrow$ plasticity (H3)}

\subsection{De-zonation proxy}
Per cell, distance of $\coordc$ from the committed ends:
\[
  \text{dez} = -\left|\frac{\coordc - \operatorname{median}(\coordc)}{\operatorname{std}(\coordc)}\right|
\]
(near the middle $\Rightarrow$ more de-zonated); high entropy is an alternative proxy.

\subsection{The confound and the within-donor test}
Both \code{dez} and \code{plasticity} rise with stage, so a \emph{pooled} correlation is a
fake-positive. Test \textbf{within donor}, then aggregate:
\begin{itemize}[leftmargin=1.4em]
  \item Per donor $d$ ($\ge 30$ cells):
        $r[d] = \operatorname{spearman}\big(\text{dez}[\text{donor}=d],\ \text{plasticity}[\text{donor}=d]\big)$;
        report $\operatorname{mean}(r)$ and the $\%$ of donors with $r > 0$.
  \item \textbf{Regression (formal version).} OLS
        $\text{plasticity} \sim \text{dez} + C(\text{stage}) + C(\text{donor})$; the \code{dez}
        coefficient is the association \emph{after} removing stage and donor; report its sign and
        $p$.
  \item \textbf{Optional two-group view.}
        $\operatorname{MannWhitneyU}(\text{plasticity}[\text{de-zonated}], \text{plasticity}[\text{zonated}])$
        within stage, with a rank-biserial effect size.
\end{itemize}

\textbf{Output A8:} a stage-stratified scatter plus the within-donor statistic and the \code{dez}
coefficient. \textbf{Acceptance:} the within-donor (not pooled) effect is reported with sign and
$p$, honestly (positive or null).

\medskip
\noindent\textbf{Code:} \code{src/steps/step8\_plasticity.py} (ref \code{src/pipeline.py:plasticity});
plot \code{src/plotting/artefacts.py:plot\_a8\_plasticity}.

\section{Step 9 --- Bonus: collapse vs inherited risk}

\subsection{Inputs}
\begin{itemize}[leftmargin=1.4em]
  \item \code{hits} $=$ the Step-7 hit list ($q < 0.05$, non-signature).
  \item \code{risk} $=$ Paper 3 risk-variant target genes.
  \item background $=$ genes \textbf{testable} in Step 7 (detected $+$ tested), ideally
        expression-matched.
\end{itemize}

\subsection{Test (hypergeometric enrichment)}
With $k = |\text{hits} \cap \text{risk}|$, $n = |\text{hits}|$,
$K = |\text{background} \cap \text{risk}|$, $N = |\text{background}|$:
\begin{align}
  p    &= \operatorname{hypergeom.sf}(k-1,\ N,\ K,\ n) && \text{(scipy)}, \\
  \text{fold} &= \frac{k/n}{K/N},
\end{align}
with BH applied across the few tests.

\textbf{Output A9:} fold $+$ CI $+$ $q$. \textbf{Acceptance:} the background is well-defined and the
result is stated with caveats (associational). Strictly optional.

\medskip
\noindent\textbf{Code:} \code{src/steps/step9\_bonus\_enrichment.py}.

\section{Plotting machinery}

\begin{center}
\begin{tabular}{@{}l l l@{}}
\toprule
\textbf{Function} & \textbf{Input} & \textbf{Draws} \\
\midrule
\code{plot\_a1\_reference\_profiles} & gene $\times$ zone means & PC rising / PP falling along zone \\
\code{plot\_a4\_coordinate}         & A4 table                & coord histogram $+$ 2 marker scatters \\
\code{plot\_a5\_validation}         & A4 healthy $+$ markers  & bar of $\rho$, colored by expected sign \\
\code{plot\_a5b\_ruler}             & per-stage diagnostics   & coherence / anti-corr / split-half curves \\
\code{plot\_a6\_collapse}           & per-donor metrics       & donor points $+$ per-stage mean \\
\code{plot\_a7\_volcano}            & DE table                & effect vs $-\log_{10} q$ volcano \\
\code{plot\_a8\_plasticity}         & A4 table                & stage-stratified dez-vs-plasticity scatter \\
\bottomrule
\end{tabular}
\end{center}

\medskip
\noindent\textbf{Code:} \code{src/plotting/artefacts.py}.

\section{Drivers}

\begin{itemize}[leftmargin=1.4em]
  \item \textbf{Full pipeline driver.} \code{src/run\_all.py} --- runs Steps 2--9 end to end and
        writes all artefacts.
  \item \textbf{Positive control.} \code{src/run\_p2\_validation.py} --- runs the ruler against
        Paper 2's own data, where the answer is known, as a sanity check on scoring and the
        classifier.
\end{itemize}

\medskip
\noindent\textbf{Code:} \code{src/run\_all.py}, \code{src/run\_p2\_validation.py}.

\end{document}
